% ============ 序言(原书序文,原书 pp.1-2 = PDF pp.3-4) ============
\chapter*{序言}\addcontentsline{toc}{chapter}{序言}

本书的目的, 是向对分析学感兴趣的读者传达调和分析学是什么. 调和分析学是以 Lebesgue 积分论为基础的理论, 不理解 Lebesgue 积分论, 就无法达到深刻的理解. 然而, 即使不使用 Lebesgue 积分论也能理解其本质的部分也不在少数, 在第1章中, 我们试图完全不使用 Lebesgue 积分来阐述调和分析. 本书所需要的 Lebesgue 积分知识, 已在第2章前半部分作了整理. 未修过 Lebesgue 积分的读者, 由于第2章前半部分正是关于 Lebesgue 积分的, 可以将其作为事实接受并继续读下去. 即便作为事实接受并继续, 我想也能在一定程度上读下去.

调和分析这个领域实际上非常广阔, 是难以把握的世界. 分析学的领域中有偏微分方程论、算子代数论、复分析、Hilbert 空间论等各种领域, 而调和分析是横跨所有这些领域的学科, 究竟以什么来称之为调和分析, 似乎并不明确. 也许可以称之为一门巧妙利用正负相消, 将由此得到的东西巧妙地累加来进行调节的学问吧?

话虽如此, 为了说明调和分析是什么, 这里原则上将"研究函数性质的数学"称为调和分析. 舍弃函数性质进行抽象化的概念之一有 Banach 空间, 这里也将其作为抽象化的调和分析纳入这一范畴. 原本创立 Banach 空间的 Banach 把一个个函数看作点, 这在当时(恐怕现在也是)是划时代的发现, 所以我认为将 Banach 空间视为调和分析是有意义的.

调和分析所研究的对象是函数的性质, 大致分为函数的微分(质)与函数的积分(量). 正如高中所学的数学那样, 微分与积分是互逆的运算, 但在分析学中这一点并不太被强调. 实际上, 微分和积分各自都有很大的延展, 将它们作为不同的东西来处理才是妥当的. 下面谈谈本书的构成. Lebesgue 积分论是研究调和分析所必需的理论, 但却是困难的理论. 因此, 本书为了不用 Lebesgue 积分也能理解, 决定使用数列空间来概观调和分析. 这就是第1章. 从第2章开始, 转入使用 Lebesgue 积分论的调和分析. 第2章将理解调和分析所必需的 Lebesgue 积分论和泛函分析学作了紧凑的整理. 如果讲解 Lebesgue 积分, 书籍就会变得庞大, 因此这里不作证明, 只记载必要事项. 关于泛函分析学, 以数列空间中触及的内容为基础, 给出了必要事项的证明. 第3章整理了调和分析的重要主题兼技巧的极大算子. 第1章考虑以全体自然数为基调的数列空间, 而第3章考虑 Euclid 空间上的 Lebesgue 空间, 因而叙述变得复杂. 也可以认为正是把那复杂的部分特意抽了出来. 第4章也可以说是考虑第1章所处理的内容在 Lebesgue 空间中会变成怎样. 数学分析中出现的积分算子以 Lebesgue 空间论中所处理的积分算子为基本, 但有时仅使用 Lebesgue 空间论中所处理的不等式, 无法得到令人满意的结论. 利用 Euclid 空间特有的性质可以进行进一步的分析, 将其应用于得到种种结果, 是很重要的事情. 其一个例子就是第1章中处理的 Calderón--Zygmund 分解. 第4章特别考察对流体力学等也很重要的 Hilbert 变换的性质. 第5章考察 Fourier 变换. Fourier 变换是数学分析中重要的积分算子, 从它的定义起就是调和分析学的重要主题. 从 Fourier 变换的定义式来看, 似乎函数不可积就无法定义, 但通过种种方法, 对不可积的函数也能定义. 此时重要的是巧妙地运用拓扑来操作极限等, 本书特别注意到这一侧面, 说明了 Fourier 变换的定义和处理方式等. 由于要处理所谓函数以外对象的 Fourier 变换, 处理会变得困难, 但本书仔细说明了应当注意哪些点来处理. 最后的第6章, 作为 Fourier 变换的应用, 说明了作为 $L^{p}$ 范数改写工具的 Littlewood--Paley 理论. 正如本书也要提到的, $L^{p}$ 范数在 $p = 1, 2$ 以外的情况下难以处理, 但对于其他的 $p$, 通过变形范数有时能看清各种各样的性质. 其一个例子就是 Littlewood--Paley 理论, 本书也兼作与第1章数列空间的比较, 说明了 Littlewood--Paley 理论.

调和分析学中有各种各样的理论, 是非常丰富的世界. 插值理论、函数空间论等作为调和分析学理论的一部分发展至今, 但考虑到入门书的性质, 除插值理论的特殊情形外, 本书不包含这些理论. 请参考卷末列举的书籍. 不过, 本书没有网罗的理论, 其背景也有与本书所说明的类似内容, 需要那种理论的动机, 可以从本书所写的内容中推测出来吧.

本书基于笔者在学習院大学、中央大学、東京都立大学的讲义写成. 执笔之际, 得到了中央大学的波多野修也先生、阿部雄斗先生, 東京都立大学的野ケ山徹先生、服部久美子先生的建议. 在此谨表谢意.

2022年1月

泽野嘉宏
